\documentclass[twocolumn]{aastex631}
\begin{document}
\title{A residual reionization ionizing-photon-budget shortfall for star-forming galaxies at z\~{}8 (crisis systematic corner)}
\author{NebulaMind Lab (autonomous pipeline)}
\affiliation{NebulaMind Lab, \url{https://lab.nebulamind.net}}
\begin{abstract}
Reionization ionizing-photon-budget reconciliation at z\~{}8: star-forming galaxies require f\_esc=0.563 (+0.449/-0.248) to close the budget (Madau-Dickinson SFRD, log xi\_ion=25.5+/-0.15, clumping C=2-5, JWST-SFRD tail), versus indirect-proxy-inferred f\_esc=0.050 (+0.075/-0.030) from LzLCS O32/beta calibrations. Median delta(required-inferred)=+0.493 dex-frac (16-84\%: +0.237 to +0.944); 99\% of the systematic MC shows a shortfall. Conclusion: a genuine SHORTFALL remains, and the sign holds under both O32 and beta calibrations. The result is bounded by the xi\_ion x clumping x proxy-calibration systematic, not by statistics. Generated autonomously from public data (jwst) via the NebulaMind Lab runner: a bounded, reproducible, descriptive study.
\end{abstract}
\section{Introduction}
Recent studies have highlighted a potential crisis in reconciling the photon budget during reionization [Muñoz2024], with increased demands on ionizing sources [Davies2021]. This has led to efforts in calibrating excursion set reionization models to conserve ionizing photons [Park2022] and assessing the galaxy ionizing photon budget at high redshifts [Duncan2015]. However, a comprehensive understanding of the reionization process remains elusive. This study aims to address this gap by examining the role of star-forming galaxies in reionization using an analytic approach [Madau2017].
\section{Data and method}
To investigate the reionization-photon-budget, we adopt a literature-anchored budget calculation without utilizing any survey catalog data. The cosmic star formation rate density (SFRD) is derived from the Madau \& Dickinson (2014) analytic fitting function. We also incorporate published values for xi\_ion and O32/beta f\_esc proxy calibrations from LzLCS [Chisholm+22, Flury+22; Simmonds+24]. Our method focuses on reconciling the ionizing-photon-budget using these parameters.
\section{Result}
Our analysis reveals that star-forming galaxies require an escape fraction of f\_esc=0.563 (+0.449/-0.248) to close the reionization photon budget at z\~{}8, assuming a Madau-Dickinson SFRD, log xi\_ion=25.5+/-0.15, and clumping factor C=2-5. In contrast, indirect-proxy-inferred f\_esc values from LzLCS O32/beta calibrations yield f\_esc=0.050 (+0.075/-0.030). The median delta between required and inferred escape fractions is +0.493 dex-frac (16-84\%: +0.237 to +0.944), with 99\% of systematic Monte Carlo simulations showing a shortfall.
\begin{figure}\centering\includegraphics[width=\columnwidth]{result.png}\caption{A residual reionization ionizing-photon-budget shortfall for star-forming galaxies at z\~{}8 (crisis systematic corner)}\end{figure}
\section{Caveats}
It is essential to acknowledge the limitations of our approach, which relies on automated single-selection and uncalibrated measurements. The accuracy of our results depends heavily on the adopted values for xi\_ion and O32/beta f\_esc proxy calibrations, as well as the clumping factor C. Additionally, our study does not account for potential variations in these parameters across different galaxy populations or redshifts. Further research is needed to refine these estimates and better understand the complexities of reionization.
\section*{References}
\footnotesize
[Muoz2024] Muñoz et al. (2024). Reionization after JWST: a photon budget crisis?. bibcode 2024MNRAS.535L..37M \par [Davies2021] Davies et al. (2021). The Predicament of Absorption-dominated Reionization: Increased Demands on Ionizing Sources. bibcode 2021ApJ...918L..35D \par [Park2022] Park et al. (2022). Calibrating excursion set reionization models to approximately conserve ionizing photons. bibcode 2022MNRAS.517..192P \par [Duncan2015] Duncan et al. (2015). Powering reionization: assessing the galaxy ionizing photon budget at z < 10. bibcode 2015MNRAS.451.2030D \par [Madau2017] Madau (2017). Cosmic Reionization after Planck and before JWST: An Analytic Approach. bibcode 2017ApJ...851...50M

\end{document}
